% БҮЛЭГ 1: СУДАЛГААНЫ ХЭСЭГ
\chapter{СУДАЛГААНЫ ХЭСЭГ}
\label{Chapter1}
\pagecolor{white}

\newcommand{\keyword}[1]{\textbf{#1}}
\newcommand{\tabhead}[1]{\textbf{#1}}
\newcommand{\code}[1]{\texttt{#1}}

%===============================================================
\section{1.1 Үндэслэл}

Хүнсний үйлчилгээний салбар нь дэлхийн нийт ДНБ-ний 3.4\%-ийг бүрдүүлж, 2023 оны
байдлаар 3.5 их наяд ам.долларын зах зээлийн хэмжээтэй байна \cite{statista2023}.
Монгол улсад тухайн салбар ДНБ-ний ойролцоогоор 2.1\%-ийг бүрдүүлж байгаа бөгөөд
Улаанбаатар хотын 5,200 гаруй хүнсний үйлчилгээний газар 45,000 гаруй ажилтанг
хамарч байна \cite{mrhba2023}.

COVID-19 цар тахлын үе нь дэлхийн рестораны салбарт гар хүрэлтгүй үйлчилгээний
эрэлтийг огцом нэмэгдүүлсэн бөгөөд QR кодоор тоног төхөөрөмж хоорондын харилцааг
хялбаршуулах боломж тодорхой болсон \cite{gossling2022}. Азийн орнуудад QR кодын
ашиглалт 2020--2023 оны хооронд 185\%-иар өссөн бөгөөд ресторан, худалдаа,
тээврийн салбарт өргөн нэвтэрч байна \cite{juniper2023}.

Гэвч Монгол улсын рестораны 92\%-д шинэ технологийг нэвтрүүлэхэд санхүүгийн
дарамт, техникийн мэдлэгийн дутагдал, Монгол хэлийг дэмждэг бэлэн шийдлийн
байхгүй байдал зэрэг саад тотгор хэвээр байна. Иймд Монголын онцлогт тохирсон,
хямд өртөгтэй, нэвтрүүлэхэд хялбар шийдэл боловсруулах нь тулгамдсан асуудал
болж байна.

%===============================================================
\section{1.2 Системийн зорилго}

Энэхүү дипломын ажлын гол зорилго нь рестораны үйлчлүүлэгч ширээн дэлгэц дээрх
QR кодыг смартфоноор уншуулахад тухайн рестораны дижитал цэс нэн даруй нээгдэж,
хоол сонгон захиалга өгөхөд тухайн захиалга нь бодит цаг хугацаанд ажилтны
дэлгэцэнд дамждаг бүрэн автоматжсан системийг загварчилж хэрэгжүүлэх явдал юм.

%===============================================================
\section{1.3 Системийн зорилт}

\begin{enumerate}
    \item Холбогдох шинжлэх ухааны бүтээлд дүн шинжилгээ хийж, санал болгах
          системийн онолын үндэслэлийг тогтоох;
    \item Хэрэглэгчийн болон техникийн шаардлагыг тодорхойлох;
    \item QR кодоос захиалга ажилтанд хүртэлх бүрэн урсгалын системийн загварыг
          боловсруулах;
    \item 4 давхаргат системийн архитектурыг хэрэгжүүлэх;
    \item Уламжлалт системтэй харьцуулан туршилтын үр дүнд дүн шинжилгээ хийх.
\end{enumerate}

%===============================================================
\section{1.4 Ижил төстэй системийн судалгаа}

\subsection{1.4.1 Монголын системүүдтэй харьцуулсан судалгаа}

Монгол улсад хоол захиалгын системийн хувьд голлон уламжлалт арга буюу цайчинд
суурилсан захиалга зонхилж байна. Хэдийгээр Facebook, Instagram зэрэг сошиал
медиагаар онлайн хүргэлтийн захиалга авдаг рестораны тоо нэмэгдсэн ч ресторан
доторх ширээнээс шууд захиалга авдаг системийн хэрэгжилт хаана ч байдаггүй.

Нямбаяр нар \cite{nyambayar2022} нарын 2022 оны судалгаагаар Улаанбаатар хотын
рестораны дижитал шилжилтийн бэлэн байдалд дүн шинжилгээ хийж, QR нэвтрүүлэх
суурь нөхцөл (смартфоны өндөр хүртээмж, 4G сүлжээний тогтвортой байдал) бүрдсэн
боловч тохиромжтой нутагшуулсан шийдэл байхгүй байгааг онцолжээ.

\subsection{1.4.2 Гадны системүүдтэй харьцуулсан судалгаа}

Дэлхийд одоо ашиглагдаж буй гол шийдлүүдийг Хүснэгт~\ref{tab:compare} дээр харуулав.

\begin{table}[htbp]
\caption{Ижил төстэй системүүдийн харьцуулалт}
\label{tab:compare}
\centering
\begin{tabularx}{\textwidth}{|>{\bfseries}X|Z|Z|Z|Z|}
\hline
\rowcolor{TableHeader}\textcolor{white}{\textbf{Шалгуур}}
  & \textcolor{white}{\textbf{Toast POS}}
  & \textcolor{white}{\textbf{Lightspeed}}
  & \textcolor{white}{\textbf{Square}}
  & \textcolor{white}{\textbf{Энэхүү систем}} \\ \hline
QR захиалга          & Нэмэлт    & Нэмэлт    & Тийм      & Гол функц \\ \hline
Бодит цаг WebSocket  & Хэсэгчлэн & Тийм      & Хэсэгчлэн & Тийм \\ \hline
Монгол хэл           & Үгүй      & Үгүй      & Үгүй      & Тийм \\ \hline
QPay / SocialPay     & Үгүй      & Үгүй      & Үгүй      & Тийм \\ \hline
Сарын зардал (USD)   & 69--165   & 69--135   & 60+       & \textasciitilde10--25 \\ \hline
Нээлттэй эх код      & Үгүй      & Үгүй      & Үгүй      & Тийм \\ \hline
\end{tabularx}
\end{table}

\textbf{Toast POS (АНУ, 2012)} нь АНУ-ын өргөн хэрэглэгддэг POS системүүдийн нэг.
QR кодоор захиалга авах боломжтой боловч нэмэлт модул шаарддаг. Монгол хэл болон
QPay, SocialPay зэрэг локал төлбөрийн хэрэгсэл байхгүй.

\textbf{Lightspeed Restaurant (Канад, 2005)} нь 115,000 гаруй рестораны газарт
хэрэглэгддэг. Сарын суурь хөлс \$69 бөгөөд QR захиалгын модул нэмэлтийн \$39
хөлстэй. Гадаад валютаар тооцогдох зардлын дарамт практикт саад учруулна.

\textbf{Square for Restaurants (АНУ, 2018)} нь QR захиалгыг дэмждэг боловч
гадаад картын системд суурилсан тул Монгол улсын картын дэд бүтцэд зориулагдаагүй.

%===============================================================
\section{1.5 Хууль, дүрэм, журам}

\subsection{1.5.1 Хувийн мэдээлэл хамгаалах тухай хууль}

Монгол улсын Хувийн мэдээлэл хамгаалах тухай хуулийн дагуу системийн хэрэглэгчийн
мэдээллийг зөвхөн зөвшөөрөлтэйгээр цуглуулж, хадгалах ёстой. Тус систем нь зочны
хувийн мэдээллийг шаардалгүй (нэвтрэлтгүй горимд) захиалга авдаг тул уг хуулийн
шаардлагыг хамгийн доод түвшинд хангах боломжтой \cite{mn_privacy2021}.

\subsection{1.5.2 Кибер аюулгүй байдлын тухай хууль}

Монгол улсын Кибер аюулгүй байдлын тухай хуулийн дагуу мэдээллийн системийн
аюулгүй байдлыг хангах, мэдээллийг зөвшөөрөлгүй хандалтаас хамгаалах үүрэг
хэрэглэгчдэд хүлээгддэг. Тус систем нь HTTPS (TLS 1.3), JWT token баталгаажуулалт,
RBAC хандалтын удирдлага ашиглан уг шаардлагыг хангана.

\subsection{1.5.3 Цахим төлбөр тооцооны тухай хууль}

Монгол улсын Цахим төлбөр тооцооны тухай хуульд заасны дагуу цахим төлбөрийн
систем нь санхүүгийн зохицуулалтын байгууллагад бүртгэгдсэн байх шаардлагатай.
Тус систем нь бүртгэгдсэн QPay болон SocialPay API-г холбоно.

\subsection{1.5.4 Хүнсний тухай хуулийн холбогдох заалт}

Хүнсний тухай хуульд заасны дагуу рестораны цэс нь хоолны найрлага, калорийн
мэдээлэл агуулсан байвал зохих. Тус систем нь цэсийн загварт найрлага, алергийн
мэдээлэл оруулах бүрэн боломжтой.

%===============================================================
\section{1.6 Онолын хэсэг}

\subsection{1.6.1 QR Кодын Онол ба Технологи}

QR код (Quick Response Code) нь 1994 онд Японы Denso Wave компанийн инженер
Масахиро Хара удирдлагаар боловсруулсан хоёр хэмжээст матрицан баркодын стандарт
юм \cite{iso18004}. Уламжлалт нэг хэмжээст баркодоос ялгаатай нь QR код нь хэвтээ
ба босоо аль аль тэнхлэгт мэдээлэл агуулдаг.

\begin{table}[htbp]
\caption{QR кодын техникийн онцлог}
\label{tab:qrspec}
\centering
\begin{tabularx}{0.85\textwidth}{|X|X|}
\hline
\rowcolor{TableHeader}\textcolor{white}{\textbf{Шинж чанар}} & \textcolor{white}{\textbf{Тодорхойлолт}} \\ \hline
Мэдээллийн хэмжээ    & Хамгийн ихдээ 4,296 тэмдэгт (ASCII) \\ \hline
Алдааны засралт      & L(7\%), M(15\%), Q(25\%), H(30\%) \\ \hline
Загварын хувилбар    & 1 (21$\times$21) -- 40 (177$\times$177) \\ \hline
Уншигч               & Смартфоны камерийн програм хангамж \\ \hline
Уншигдах өнцөг       & Бүх чиглэл (360°) \\ \hline
\end{tabularx}
\end{table}

Алдааны засралтын H түвшинд кодын 30\% нь бүрхэгдсэн ч мэдээллийг бүрэн уншиж
чадна. Рестораны ширээн дэлгэцэнд ашиглахад гэмтэх, бохирдох эрсдэл бага тул
M (15\%) түвшин оновчтой.

\subsection{1.6.2 WebSocket Протоколын Онол}

WebSocket нь Fette \& Melnikov \cite{rfc6455} нарын RFC 6455 стандартаар
тодорхойлогдсон, клиент болон сервер хооронд нэг TCP холболтоор хоёр чиглэлт
(bidirectional) тасралтгүй харилцааг боломжтой болгодог протокол юм.

\begin{table}[htbp]
\caption{HTTP ба WebSocket харьцуулалт}
\label{tab:wscompare}
\centering
\begin{tabularx}{\textwidth}{|X|Z|Z|}
\hline
\rowcolor{TableHeader}\textcolor{white}{\textbf{Шинж чанар}}
  & \textcolor{white}{\textbf{HTTP}}
  & \textcolor{white}{\textbf{WebSocket}} \\ \hline
Холболтын чиглэл      & Нэг чиглэлт    & Хоёр чиглэлт \\ \hline
Тогтвортой холболт    & Хүсэлт бүр шинэ & Тасралтгүй \\ \hline
Header overhead       & 200--800 байт  & Зөвхөн эхлэлд \\ \hline
Бодит цагийн дэмжлэг & Polling шаардлагатай & Байгалийн \\ \hline
Хоцролт (latency)     & 50--300 ms     & 1--50 ms \\ \hline
\end{tabularx}
\end{table}

Hu нар \cite{hu2023} 100 зэрэг клиентийн нөхцөлд WebSocket нь HTTP long-polling-тай
харьцуулахад серверийн CPU ашиглалт 52\% бага, хариу хугацаа 3.2 дахин хурдан
байсан болохыг туршлагаараа нотолжээ.

\subsection{1.6.3 RESTful API Архитектурын Онол}

REST (Representational State Transfer) нь Roy Fielding \cite{fielding2000} нарын
докторын зэрэгт зориулсан бүтээлд тодорхойлсон вэб архитектурын хэв маяг юм.
REST нь 6 архитектурын хязгаарлалтад суурилна: Client-Server, Stateless, Cacheable,
Uniform Interface, Layered System, Code on Demand.

\subsection{1.6.4 JWT (JSON Web Token) Нэвтрэлтийн Технологи}

JWT нь RFC 7519 \cite{rfc7519} стандартаар тодорхойлогдсон, JSON форматад мэдээлэл
агуулсан баталгаажуулалтын token юм. Гурван хэсгийг цэгээр тусгасан: Header (алгоритм),
Payload (мэдээлэл), Signature (гарын үсэг). JWT-ийн \textbf{stateless} байдал нь
мобайл аппликейшнд session cookie-г ашиглахгүйгээр нэвтрэлтийг хэрэгжүүлэх
боломж олгодог.

%===============================================================
\section{1.7 Технологийн судалгаа}

\begin{table}[htbp]
\caption{Ашигласан технологиудын жагсаалт}
\label{tab:techstack}
\centering
\begin{tabularx}{\textwidth}{|>{\bfseries}X|X|X|}
\hline
\rowcolor{TableHeader}\textcolor{white}{\textbf{Технологи}}
  & \textcolor{white}{\textbf{Хувилбар}}
  & \textcolor{white}{\textbf{Зориулалт}} \\ \hline
React Native + Expo & 0.73 / SDK 50 & Мобайл клиент (iOS + Android) \\ \hline
Node.js + Express   & 20 LTS / 5.x  & API сервер \\ \hline
Socket.IO           & 4.x           & WebSocket бодит цагийн харилцаа \\ \hline
PostgreSQL          & 16            & Харилцаат өгөгдлийн сан \\ \hline
Drizzle ORM         & 0.30          & Type-safe өгөгдлийн сангийн давхарга \\ \hline
JWT (jsonwebtoken)  & 9.x           & Хэрэглэгчийн нэвтрэлт \\ \hline
Zod                 & 3.x           & Input validation \\ \hline
\end{tabularx}
\end{table}

\textbf{React Native} нь Meta (Facebook) компанийн 2015 оноос нийтэд нээсэн
JavaScript дээр суурилсан cross-platform мобайл хөгжүүлэлтийн framework юм.
"Learn once, write anywhere" зарчмаар iOS болон Android аль алинд нэг кодын
санд хөгжүүлэх боломжтой.

\textbf{Node.js} нь Chrome V8 JavaScript engine дээр суурилсан серверийн runtime.
Event-driven, non-blocking I/O загвар нь олон зэрэг хүсэлтийг хөнгөн арга хэлбэрээр
зохицуулдаг тул рестораны систем шиг олон зэрэг хэрэглэгчийн нөхцөлд тохиромжтой.

\textbf{PostgreSQL} нь 1996 оноос хойш хөгжүүлэгдэж буй нээлттэй эх кодтой
object-relational өгөгдлийн сан бөгөөд ACID баталгааг бүрэн хэрэгжүүлдэг.
Массив төрлийн өгөгдөл (алергийн жагсаалт), JSON дэмжлэг тус системд тохиромжтой.

%===============================================================
\section{1.8 Бүлгийн дүгнэлт}

Энэхүү бүлэгт судалгааны үндэслэл, зорилго, зорилтуудыг тодорхойлж, олон улсын
болон Монголын ижил төстэй системүүдтэй харьцуулсан шинжилгээ хийв. QR код,
WebSocket, RESTful API болон JWT-ийн онолын суурийг тодорхойлов. Харьцуулалтаар
Toast POS, Lightspeed, Square зэрэг гадаадын шийдлүүд Монголын онцлогт
(хэл, локал төлбөр, зардал) нийцэхгүй байгаа нь нотлогдсон тул нутагшуулсан
шийдэл боловсруулах шаардлага тодорхой болов.
